% !TeX program = xelatex
% UTF-8; compile twice with XeLaTeX. Requires ctex and Fandol fonts.
\documentclass[UTF8,a4paper,zihao=-4,fontset=fandol]{ctexart}
\usepackage[margin=2.3cm]{geometry}
\usepackage{enumitem,booktabs,tabularx,longtable,array,fancyhdr}
\usepackage[hidelinks]{hyperref}
\setlength{\parindent}{2em}
\setlength{\parskip}{0.35em}
\setlist[enumerate]{leftmargin=2em,itemsep=0.35em,topsep=0.3em}
\setlist[itemize]{leftmargin=2em,itemsep=0.25em,topsep=0.3em}
\renewcommand{\arraystretch}{1.3}
\pagestyle{fancy}
\fancyhf{}
\fancyhead[L]{\small 星知智链｜作品 Demo}
\fancyhead[R]{\small 草稿・待现场验收}
\fancyfoot[C]{\thepage}
\setlength{\headheight}{15pt}
\newcommand{\pending}[1]{\textbf{【待补：#1】}}
\begin{document}
\begin{center}
  {\LARGE\bfseries 03—作品 Demo}\par\vspace{0.7em}
  {\large 星知智链——多源知识融合与智能推理平台}\par
  可访问入口、双角色操作说明与演示验收草稿
\end{center}
\noindent\begin{tabularx}{\textwidth}{@{}lX@{}}
\toprule
赛事 & 2026年度中国青年科技创新“揭榜挂帅”擂台赛（学生赛道） \\
作品编码／赛题 & 592374 ／ XH-202620 \\
推报学校／申报者 & 武汉大学 ／ 赵晨旭 \\
体验入口 & \url{https://esa.lovelearnlearning.cn/} \\
编制日期 & 2026年9月11日 \\
验收状态 & 仓库实现已核对；当前站点、账号及模型全流程待实测 \\
\bottomrule
\end{tabularx}

\section{作品定位与演示目标}
星知智链（ESA）是面向计算机学科的教学、学习与科研辅助平台。系统将课程知识图谱、学习证据、文档检索和智能体工具结合起来，在教师、学生及项目资源边界内提供辅助服务。

本次 Demo 以\textbf{“作业诊断—教师复核—反馈发布—针对性学习”}为主线，重点展示：AI 建议能够被教师修改；未经发布的分析不会成为学生正式成绩；正式反馈可形成学习证据，支持学生继续学习与教师查看班级知识点诊断。演示不以自动评分取代教师责任，也不将系统估计等同于实际学习成效。

\section{访问方式与账号准备}
建议使用桌面浏览器，并用两个独立浏览器会话分别登录教师和学生账号，避免角色切换时复用登录状态。该方式是演示组织建议，不代表已完成跨浏览器兼容性认证。

\noindent\begin{tabularx}{\textwidth}{@{}p{2.2cm}Xp{3.2cm}@{}}
\toprule
角色 & 用途与可访问范围 & 待填写信息 \\
\midrule
教师临时账号 & 本人演示班级、作业、AI 分析、复核发布及班级看板 & \pending{实际用户名} \\
学生临时账号 & 接受班级邀请、本人作答、已发布反馈及学习空间 & \pending{实际用户名} \\
\bottomrule
\end{tabularx}

\begin{itemize}
  \item 两个账号应仅包含合成演示数据，不提供真实学生或科研用户数据。账号尚未由本草稿创建。
  \item 密码通过赛事允许的独立安全渠道交付，不写入本文件、公开仓库或视频；不要要求评委依赖现场邮箱验证码才能进入预设体验流程。
  \item 体验有效期：\pending{起止日期与时间}；维护联系人及反馈渠道：\pending{团队确认的交付渠道}。
  \item 初审或演示结束后按约定重置密码或停用账号，并处理演示数据。
\end{itemize}

\section{演示数据与起始状态}
建议准备一个虚拟班级“数据结构 Demo 班”，关联部署环境中实际存在的数据结构课程；准备一份“二分查找诊断”作业。采用团队自编题目和虚拟作答，不依赖真实学生信息。

\noindent\begin{tabularx}{\textwidth}{@{}p{2.4cm}X@{}}
\toprule
项目 & 建议内容 \\
\midrule
题型与满分 & 简答题，10分。 \\
题目 & 在有序数组中进行二分查找，为什么最坏情况下的时间复杂度是 $O(\log n)$？请说明搜索区间如何变化。 \\
教师参考答案 & 每轮比较后，剩余搜索区间至多约为原来的一半。进行 $k$ 轮后规模约为 $n/2^k$；当其不超过1时终止，因此轮数为 $O(\log n)$。 \\
教师评分参考 & 区间减半4分；给出规模递推或轮数关系4分；复杂度结论及表达2分。此为人工演示评分参考，不是实测模型评分。 \\
虚拟学生作答 & “每次都从中间找，所以查找很快，时间复杂度是 $O(\log n)$。” \\
关联知识点 & 在实际知识图谱中选择“二分查找”或对应有效知识点。不得直接假定测试代码中的标识适用于线上索引。 \\
\bottomrule
\end{tabularx}

为方便重复体验，可预置一份\textbf{已发布且已复核的样例}供快速查看，另准备一份未完成的作业进行现场提交、分析与发布。预置样例应明确标识，不冒充当次实时生成结果。重复演示应使用新的作业或提交，不依赖尚未实现的完整反馈更正前端流程。

\section{推荐操作流程}
以下为预期交互及验收点，不是本次已完成的线上运行记录。
\begin{enumerate}
  \item \textbf{教师创建班级与邀请。}教师登录后进入“教学空间 → 教学工作台”，创建演示班级，输入学生临时账号的精确用户名并发起邀请。验收点：邀请仅面向指定账号，不使用公开邀请码扩大访问范围。
  \item \textbf{学生确认加入。}学生进入“学习空间 → 作业”，接受邀请。验收点：学生可查看已加入班级的开放作业，不能进入教师管理接口。
  \item \textbf{教师发布诊断作业。}创建上述简答题，填写满分、评分标准、参考答案与有效知识点后发布。验收点：学生可看到题目，但看不到教师参考答案及评分规则字段。
  \item \textbf{学生提交答案。}学生打开作业，填写虚拟作答并提交。验收点：产生可回看的本人提交，其他学生不能读取该提交。
  \item \textbf{教师启动 AI 分析。}教师进入作业批改页，对该提交或当前作业最新提交批量分析。验收点：出现逐题建议分、反馈等结构化结果；现场核对来源是辅助模型还是确定性降级，不能仅凭界面出现分数认定模型调用成功。
  \item \textbf{教师复核与修改。}点击“复核并调整”，根据答案修改分数、反馈和知识点，确认复核。建议反馈可写为：“请补充每轮区间减半与 $n/2^k$ 的关系。”不预先规定模型分数，也不要求人工最终分数与模型一致。
  \item \textbf{教师主动发布反馈。}发布前，学生端不应显示未发布的正式成绩或内部分析；发布后，学生可查看教师确认的逐题反馈。验收点：复核与发布是独立动作，AI 不自行发布成绩。
  \item \textbf{学生继续针对性学习。}学生查看反馈后，进入学习助手或知识地图，围绕“为什么区间减半对应对数复杂度”继续提问。验收点：正式作业证据可用于个人掌握度更新；不要求单条证据必然引起可见的等级变化，也不把任意聊天自动当成掌握证据。
  \item \textbf{教师查看班级诊断。}返回班级看板，查看已发布反馈形成的知识点聚合。前置知识点没有直接证据时应标为待诊断；单个演示学生的数据不用于推断班级整体效果或教学因果关系。
\end{enumerate}

\section{可选补充体验与功能边界}
\noindent\begin{tabularx}{\textwidth}{@{}p{2.5cm}Xp{5.1cm}@{}}
\toprule
能力 & 可展示内容 & 启用条件及限制 \\
\midrule
学习对话 & 多轮流式回答、数学公式、代码及知识地图 & 需实际模型服务可用；页面显示不等同于推理成功。 \\
文档与检索 & 上传授权材料、来源定位、公共或个人知识库检索 & 依赖相关开关、解析服务与索引；无检索证据时不能宣称已引用知识库。 \\
科研辅助 & 项目组织、前沿追踪、写作及数据分析任务 & 依赖项目权限和相应服务；不将待补引用写成真实文献。 \\
代码工具 & 代码编辑、按环境提供的补全与受控执行 & LSP 按语言服务探测；沙箱默认关闭，代码文本题不等同于自动运行判题。 \\
\bottomrule
\end{tabularx}

当前教学主线以简答题和代码文本题为主，不将多教师协作、助教角色、选择题、附件题、自动运行评分或大班级性能验收列为已完成能力。RAG、多模态、个人知识库与沙箱是否启用，以实际部署和当次验收为准，不直接沿用示例配置作结论。

\section{三分钟录屏建议}
以下为拟录制时间安排，成片长度与界面结果以真实录制为准。建班、邀请和样例准备可提前完成并说明。

\noindent\begin{tabularx}{\textwidth}{@{}p{2.5cm}X@{}}
\toprule
时间 & 画面与讲述重点 \\
\midrule
00:00--00:20 & 显示作品名称与真实入口，说明教师、学生双角色及预置的合成班级。 \\
00:20--00:45 & 学生打开诊断作业并提交，展示待处理状态。 \\
00:45--01:20 & 教师启动分析，显示真实结果与来源；说明 AI 仅给出建议。 \\
01:20--01:55 & 教师修改反馈与分数，确认复核并主动发布，突出人工裁决。 \\
01:55--02:30 & 学生查看已发布反馈，再进入学习助手或知识地图继续学习。 \\
02:30--03:00 & 教师查看班级诊断，说明证据来源、演示数据与功能边界。 \\
\bottomrule
\end{tabularx}

若模型等待时间较长，可标明等待时长并剪去等待段，但不能拼接出不存在的成功路径；出现确定性降级时应如实说明。视频不得出现凭证、真实学生信息或后台敏感日志。

\section{提交前验收记录（待执行）}
仓库材料记载：2026年9月8日首页及 \path{/api/health} 曾返回 HTTP 200。本草稿未重新检查站点，也未进行真实双账号端到端验收；历史记录不代表2026年9月11日或评审期间服务持续可用。健康检查仅证明 HTTP 进程存活，不检查模型、Qdrant、MinerU、邮件或完整教学流程。

\noindent\begin{longtable}{@{}p{3.2cm}p{9cm}p{2cm}@{}}
\toprule
检查项 & 通过标准 & 结果 \\
\midrule
\endhead
入口与登录 & HTTPS 页面可加载；教师、学生临时账号均可登录。 & 待验收 \\
对话与模型 & 实际发送请求可获得流式回答；记录服务版本与失败情况。 & 待验收 \\
教学闭环 & 邀请、发布、提交、分析、复核、反馈发布及学生查看均可完成。 & 待验收 \\
人工控制 & 未发布结果对学生不可见；教师主动发布后正式反馈才可见。 & 待验收 \\
权限隔离 & 其他教师及其他学生不能读取不属于其权限范围的提交。 & 待验收 \\
学习证据 & 发布反馈形成知识点证据；重复发布不重复计入。 & 待验收 \\
条件性能力 & 分别记录 RAG、个人库、多模态、LSP、沙箱的启用及实测状态。 & 待验收 \\
隐私与交付 & 样例无真实隐私；凭证分开交付；视频无敏感信息。 & 待验收 \\
\bottomrule
\end{longtable}
\noindent 验收人：\underline{\hspace{3cm}}\quad 时间：\underline{\hspace{4cm}}\par
\noindent 部署版本与证据位置：\underline{\hspace{10cm}}

\section{故障与降级说明}
\begin{itemize}
  \item \textbf{页面或登录失败：}核对域名、TLS、静态资源、账号及会话状态；向维护联系人反馈，不公开密码。备用录屏只能作为标明时间的历史演示，不能替代在线可用性结论。
  \item \textbf{模型超时或 AI 分析降级：}记录实际错误；确定性降级通常提供低置信度建议，教师仍需人工复核。不能将关键词匹配结果计为大模型成功分析。
  \item \textbf{检索无来源：}说明本次未取得可核验检索证据。通用知识解释不冒充已检索内容；必要时停止给出引用。
  \item \textbf{代码执行不可用：}保留文本编辑或可用的补全功能，明确代码未在隔离环境运行，不宣称已经执行或通过测试。
  \item \textbf{教师发现反馈有误：}停止继续传播错误结论，由教师组织更正；当前完整反馈更正前端流程不在已完成范围，不承诺一键撤回或无痕覆盖。
\end{itemize}

\appendix
\section{实现依据与复现提示（内部核对用）}
本草稿以工作区提交 \texttt{0c1da87} 为代码核对基线，身份信息取自用户提供的报名表。原有 Markdown／DOCX 未被本草稿覆盖，正式提交前应统一各格式版本。
\begin{itemize}
  \item \path{TEACHING_STUDENT_DEMO.md}：界面入口、角色、数据模型、接口与已知限制。
  \item \path{backend/tests/test_teaching_workflow.py}：双角色闭环、隐藏字段、发布反馈及其他教师访问限制测试；其中样例分数来自测试环境，不能挪作线上模型效果。
  \item \path{backend/core/services/teaching_analysis_service.py}：\texttt{auxiliary\_llm} 与 \texttt{deterministic\_fallback} 两种分析来源。
  \item \path{README.md}、\path{REQUEST.md}：运行配置、健康检查边界及比赛交付要求。
\end{itemize}
可在依赖齐备的隔离测试环境中运行下列命令核对教学后端流程；该测试不替代真实浏览器、模型与部署验收：
\begin{quote}\small\ttfamily
python -m pytest backend/tests/test\_teaching\_workflow.py -q
\end{quote}
部署复现遵循仓库 \path{README.md} 和 \path{deploy/} 中的配置说明；不在 Demo 文档中放入生产配置、模型 Token 或真实数据库。
\end{document}
